% Capítulo 2 — Contexto (orçamento: ~20 pp)
% O que já existe, e porque é que não chega. A matriz das três perguntas é o instrumento central.
% Regra de escrita: sem travessões em prosa, decimais com ponto, PT-PT.

\chapter{Estado da arte}
\label{cap:contexto}

Este capítulo responde a uma pergunta simples de enunciar e difícil de responder com honestidade:
\emph{o que é que já existe, e porque é que não chega?} A resposta organiza-se em três movimentos.
Primeiro caracteriza-se a pessoa a quem o sistema se dirige, porque as limitações dela são o que
define o problema. Depois examinam-se as ferramentas que hoje ocupam este espaço, com um critério
fixo em vez de uma descrição geral. Por fim percorrem-se, uma a uma, as áreas técnicas de onde vêm
as peças que o Capítulo~\ref{cap:metodologia} vai usar, para que fique claro que nenhuma delas é
invenção deste trabalho e que a contribuição está na escolha, na integração e na avaliação crítica.

\section{Quem investe, e porque é que a atenção é o problema}
\label{sec:ctx_investidor}

A escala foi dada no Capítulo~\ref{cap:introducao}: a maior parte dos adultos americanos detém
ações \autocite{gallup2025stock}, num mercado que vale mais de 62 biliões de dólares
\autocite{sifma2025factbook}. Interessa agora a segunda metade da caracterização, que não é de
dimensão mas de comportamento, e que explica porque é que o problema deste trabalho existe.

Este público é numeroso e é ativo. Mesmo durante a queda de março de 2020, os investidores
particulares de uma plataforma sem comissões aumentaram as suas posições agregadas em vez de fugirem
\autocite{welch2022robinhood}, o que contraria a caricatura do investidor particular que só reage em
pânico. São pessoas que decidem, com regularidade, e que por isso beneficiam de melhor informação.

O que a literatura de finanças comportamentais estabelece é que essas decisões se afastam de forma
sistemática, e não aleatória, do ideal racional \autocite{barberis2003behavioral}. Duas dessas
regularidades importam diretamente para o desenho de um sistema de alertas.

A primeira é a \textbf{aversão à perda}: sentimos as perdas com mais intensidade do que os ganhos
equivalentes \autocite{kahneman1979prospect}. Daí que uma pessoa fique particularmente sensível a
más notícias, e daí que um alerta que chegue a tempo e venha explicado tenha valor real em vez de
valor decorativo.

A segunda é mais desconfortável, e é a que fundamenta este trabalho inteiro. \textcite{barber2008glitters}
mostram que os investidores particulares são compradores líquidos de
ações que \emph{chamam a atenção}, entendendo-se por isso as que aparecem nas notícias, as que
apresentam volume invulgar e as que registam retornos extremos num único dia. O mecanismo que os
autores propõem é o do problema de procura: perante milhares de escolhas possíveis, a pessoa não
avalia todas, e recai sobre aquilo que lhe saltou à vista. A base empírica desta afirmação é
concreta: são registos de corretoras dos Estados Unidos entre 1991 e 1999, anteriores à negociação por
telemóvel e sem comissões, e o mecanismo que descrevem é o que aqui se assume transferível, não os
valores concretos. A atenção é de tal modo central neste
processo que chega a ser mensurável através do volume de pesquisas na internet
\autocite{da2011attention}.

Repare-se no que isto significa quando se junta às três perguntas do Capítulo~\ref{cap:introducao}.
Os três sinais que capturam a atenção do investidor são a notícia, o volume invulgar e o movimento
extremo do dia, e este sistema calcula os três. A diferença está no que acontece a seguir.

Hoje esses sinais chegam à pessoa em bruto, e o que a literatura documenta é que, em agregado,
estes investidores compram mais do que vendem aquilo que lhes chamou a atenção. Os autores são
explícitos quanto ao alcance da afirmação: nem toda a gente compra tudo o que lhe salta à vista, e
o que se mede é um desequilíbrio entre compras e vendas. São também explícitos quanto ao que daí
resulta, e é essa a metade que justifica este trabalho: esse padrão de compra \emph{não} gera
retornos superiores.

O que este trabalho tenta é interpor, entre o sinal e a decisão, o contexto que permite julgá-lo:
se aquilo é invulgar para aquela empresa, se foi a empresa ou o mercado, e o que se seguiu em casos
parecidos. A Figura~\ref{fig:ctx_atencao} resume esta cadeia.

Convém ainda ser exato quanto ao estatuto de cada um dos três sinais, porque não é o mesmo. A
notícia e o movimento extremo são os dois \emph{gatilhos}, ou seja podem por si levantar um alerta;
o volume é calculado e mostrado como contexto, e nunca levanta nada. A razão está medida e é dita na
Secção~\ref{sec:av_alternativas}: combinar o volume com os outros sinais só melhorou a triagem em um
de três orçamentos, e uma capacidade que ganha uma vez em três não justifica mais uma via de
interrupção.

Uma nota de atribuição sobre essa combinação, porque a ideia não é minha. Veio de um painel de
informação em tempo real que reúne sinais de origens diferentes num só indicador \autocite{worldmonitor2026}, sugerido pelo coorientador deste trabalho. Era plausível o suficiente
para ser \emph{medida} em vez de descartada por argumento, e é assim que entra nesta dissertação:
como uma alternativa que foi testada e não ficou.

\begin{figure}[ht]
\centering
\begin{tikzpicture}[
  font=\small, >={Stealth[length=2.4mm]},
  s/.style={draw, rounded corners, align=center, text width=2.55cm, minimum height=1.5cm,
            inner sep=4pt, fill=black!3},
  d/.style={draw, rounded corners, align=center, text width=2.55cm, minimum height=1.5cm,
            inner sep=4pt, fill=black!8, very thick},
]
  \node[s] (n)  {\textbf{Notícia}};
  % A caixa do volume é tracejada de propósito: o sinal é calculado e mostrado, mas não
  % levanta alertas. Desenhá-la igual às outras afirmaria três gatilhos onde há dois.
  \node[s, dashed, below=0.35cm of n] (v) {\textbf{Volume}\\\footnotesize invulgar};
  \node[s, below=0.35cm of v] (m) {\textbf{Movimento}\\\footnotesize extremo do dia};

  \node[d, right=2.1cm of v] (ctx) {\textbf{Contexto}\\\footnotesize as três perguntas};
  \node[s, right=1.9cm of ctx] (dec) {\textbf{Decisão}\\\footnotesize da pessoa};

  \draw[->, thick, black!45] (n) -- (ctx);
  \draw[->, thick, black!45] (v) -- (ctx);
  \draw[->, thick, black!45] (m) -- (ctx);
  \draw[->, thick, black!45] (ctx) -- (dec);

  \draw[->, dashed, black!55] (n.east) to[out=20, in=150] node[above, midway, font=\scriptsize,
        align=center, yshift=1pt] {hoje: sem intermediário} (dec.north west);
\end{tikzpicture}
\caption[Da atenção à decisão]{Os três indicadores de atenção que a literatura usa
\autocite{barber2008glitters} são calculados por este sistema, mas apenas dois levantam alertas: a
caixa do volume está a tracejado porque o sinal é mostrado como contexto e não interrompe ninguém. A
seta a tracejado é o caminho que existe hoje, em que o sinal chega em bruto. A caixa central é o que
este trabalho acrescenta.}
\label{fig:ctx_atencao}
\end{figure}

Uma última nota de enquadramento. Do lado das \emph{instituições}, a adoção de inteligência
artificial em serviços financeiros deixou de ser experimental e passou a ser generalizada, como
documenta o inquérito a empresas do setor conduzido pelo \textcite{ccaf2026aifs}. Do lado dos
\emph{produtos de retalho} não cito esse inquérito, porque ele não mede essa população; o que se
sabe é o que os próprios fornecedores anunciam, e disso trata a Secção~\ref{sec:ctx_produtos}. Com
essa ressalva, a pergunta relevante já não é se estas técnicas vão chegar ao investidor particular,
mas em que condições de transparência é que chegam.

\section{O que existe hoje para quem investe pouco}
\label{sec:ctx_ferramentas}

Resolveram, em parte, e de maneiras diferentes. O que existe merece ser olhado com um critério fixo,
em vez de o descrever em geral. O critério deste trabalho são as três perguntas do
Capítulo~\ref{cap:introducao}, e a Tabela~\ref{tab:ctx_matriz} usa-as para pontuar as alternativas.

\begin{table}[ht]
\caption[As ferramentas existentes contra as três perguntas]{O que cada tipo de ferramenta responde.
``Parcial'' significa que dá matéria-prima para a resposta, mas deixa o trabalho ao utilizador.}
\label{tab:ctx_matriz}
\centering
\small
\begin{tabular}{L{0.21\textwidth} C{0.11\textwidth} C{0.15\textwidth} C{0.13\textwidth}
                L{0.16\textwidth}}
\toprule
\tabhead{Ferramenta} & \tabhead{É invulgar?} & \tabhead{Empresa ou mercado?} & \tabhead{Já aconteceu?}
& \tabhead{Custo} \\
\midrule
Aplicação da corretora & não & não & não & incluído \\
Agregador de notícias & não & não & não & grátis \\
Aplicação de sentimento & parcial & não & não & grátis/baixo \\
\emph{Robo-advisor} & não & não & não & \% da carteira \\
Terminal profissional & sim & sim & sim & milhares/ano \\
\midrule
\textbf{Este trabalho} & \textbf{sim} & \textbf{sim} & \textbf{sim} & \textbf{grátis} \\
\bottomrule
\end{tabular}
\end{table}

A tabela mostra o essencial: \textbf{a linha que responde às três perguntas existe, e é a mais cara
de todas}. Não há aqui um problema científico por resolver: há um problema de acesso. As
respostas são conhecidas; estão atrás de uma subscrição que a pessoa de quem este trabalho fala não
vai pagar.

Três linhas merecem atenção em particular, porque cada uma falha por uma razão diferente e essas
razões informam decisões de desenho tomadas mais à frente.

A \textbf{aplicação da corretora} é onde a maioria das pessoas realmente olha. Mostra a percentagem do dia
e um gráfico. Não diz se aquilo é muito para aquela ação, não separa a empresa do mercado, e não tem
memória. A pergunta ``isto é normal?'' fica sem resposta precisamente no sítio onde é feita. Os
alertas de preço que estas aplicações oferecem partilham a mesma limitação de forma mais aguda:
disparam quando um limiar fixo é atravessado, o que ignora a volatilidade própria de cada ação, e
por isso disparam constantemente nas empresas agitadas e quase nunca nas calmas. Além disso dizem
apenas \emph{que} um nível foi ultrapassado, nunca \emph{porquê}. Esta observação não é acessória: é
a razão pela qual o Capítulo~\ref{cap:metodologia} usa uma medida relativa à norma de cada empresa em
vez de um limiar absoluto, e o Capítulo~\ref{cap:avaliacao} mede a diferença que isso faz.

As \textbf{aplicações de sentimento} pontuam notícias como positivas ou negativas e mostram o resultado.
Dão matéria-prima para a primeira pergunta, mas param aí, e o Capítulo~\ref{cap:avaliacao} mostra,
com medição própria, que o tema de uma notícia e a direção do movimento que se lhe segue são coisas
distintas com mais frequência do que a intuição sugere.

Os \emph{robo-advisors} \autocite{dacunto2019robo} respondem sobretudo a outra pergunta: como
distribuir dinheiro por uma carteira. É uma pergunta legítima e não é esta. A literatura
reconhece-lhes benefícios genuínos, embora não uniformes: \textcite{dacunto2019robo} medem uma
melhoria da diversificação e uma redução da volatilidade \emph{para quem tinha menos de cinco ações}
antes da adesão, ao passo que para quem já detinha mais de dez o efeito sobre a diversificação é
quase nulo. A contenção de enviesamentos comportamentais é o benefício mais consistente
\autocite{dacunto2019robo, cardillo2024robo}.

A diferença face a este trabalho é de \emph{objeto} e não de transparência. Um \emph{robo-advisor}
recomenda uma carteira e executa-a, ainda que a decisão final continue a ser autorizada pelo
investidor; aqui não se recomenda nada, e o que se entrega é a evidência sobre um acontecimento
concreto. Acresce que aconselhar sobre investimentos tem obrigações regulatórias que este trabalho
evita por desenho, como se detalha na Secção~\ref{sec:ctx_escolhas}.

\section{Os produtos que fazem exatamente esta pergunta}
\label{sec:ctx_produtos}

Existem, são recentes, e seria desonesto não os nomear.

O \textbf{Robinhood Cortex}, anunciado em março de 2025, declara na página do próprio fornecedor que
serve para \emph{``responder à velha questão de porque é que esta ação está a subir ou a descer
hoje''} \autocite{robinhood2025cortex}. Os \textbf{momentos-chave} do Google Finance, de junho de
2026, propõem-se \emph{``explicar porque é que uma ação se moveu''} \autocite{google2026finance}.

É a mesma pergunta deste trabalho, feita por empresas com muito mais recursos e muito mais
utilizadores. Vale a pena dizê-lo com clareza, em vez de fingir que este território estava vazio.

A diferença está no que é entregue. Um resumo gerado é uma \textbf{afirmação}: o utilizador lê e
acredita, ou não. O que este trabalho entrega é uma afirmação \textbf{com a evidência anexada}: os
casos passados, um a um, com as datas, as semelhanças e o que o preço fez a seguir a cada um. A
diferença não é de fluência: é de \emph{verificabilidade}.

Uma ressalva de método, porque é fácil escorregar aqui: destes produtos só se afirma o que está
escrito nas páginas dos próprios fornecedores, consultadas em data registada. Não se testaram, e por
isso não se diz o que fazem por dentro; diz-se apenas o que declaram fazer.

\section{E um assistente de IA genérico?}
\label{sec:ctx_llm}

É a objeção mais razoável de todas, e merece resposta séria.

Um assistente genérico responde bem à pergunta, no sentido em que produz um texto fluente e
plausível. O problema é o que está por trás desse texto. Ele não tem:

\begin{itemize}
  \item a \textbf{volatilidade recente} daquela ação, que é o que permite dizer se 4\% é muito;
  \item a \textbf{separação} entre o que foi o mercado e o que foi a empresa;
  \item uma \textbf{base de casos passados} para \emph{recuperar}. No melhor dos casos
        \emph{recorda} casos do treino, o que não é a mesma coisa e não é verificável.
\end{itemize}

E há uma diferença mais funda, que é a razão pela qual este trabalho não é um invólucro à volta de um
modelo de linguagem. Modelos grandes especializados em finanças existem
\autocite{wu2023bloomberggpt} e são impressionantes. Mas a saída deles é uma afirmação que se lê bem
e não se consegue conferir, e a literatura da área avisa que uma explicação que não se pode verificar
não resolve o problema que diz resolver \autocite{lipton2018mythos, rudin2019stop}. A literatura
sobre confiança em sistemas automáticos acrescenta a razão pela qual isso importa.
\textcite{lee2004trust} descrevem \emph{dois} modos de falha simétricos: confiar de menos, e por isso
ignorar um aviso correto, e confiar de mais, e por isso aceitar um errado. O objetivo não é maximizar
a confiança nem minimizá-la, mas torná-la \emph{apropriada}, ou seja ajustada à capacidade real do
sistema. E uma pessoa só pode ajustar a sua confiança se lhe for mostrado em que é que o sistema se
apoia.

Por isso, neste trabalho, a linguagem nunca produz factos. Os números vêm dos motores que os
calcularam, e o texto é montado a partir deles, como se descreve na
Secção~\ref{sec:sis_explicacao}.

\section{Detetar o que é invulgar}
\label{sec:ctx_anomalias}

A deteção de anomalias é uma área estabelecida, com uma taxonomia consolidada que separa os métodos
pela forma como definem normalidade \autocite{chandola2009anomaly, pang2021deep}. Para o problema
deste trabalho interessam três famílias, resumidas na Tabela~\ref{tab:ctx_anomalias}.

\begin{table}[ht]
\caption[Famílias de deteção de anomalias]{As três famílias relevantes para séries temporais
financeiras, e o que cada uma exige em troca do que oferece.}
\label{tab:ctx_anomalias}
\centering
\small
\begin{tabular}{L{0.20\textwidth} L{0.30\textwidth} L{0.36\textwidth}}
\toprule
\tabhead{Família} & \tabhead{Como define normalidade} & \tabhead{Custo e benefício} \\
\midrule
Estatística \newline (limiar sobre a norma recente)
  & Distribuição local dos retornos, estimada numa janela deslizante
  & A explicação é o próprio cálculo; supõe uma distribuição localmente estável \\
\addlinespace
Aprendida, não supervisionada \newline \autocite{liu2008isolation, breunig2000lof}
  & Estrutura geométrica dos dados: pontos fáceis de isolar ou pouco densos face aos vizinhos
  & Não precisa de rótulos e capta padrões multivariados; a decisão deixa de ser legível em uma frase \\
\addlinespace
Profunda \newline \autocite{pang2021deep}
  & Um modelo de normalidade aprendido por uma rede
  & Maior capacidade; exige mais dados e afasta-se do requisito de explicabilidade \\
\bottomrule
\end{tabular}
\end{table}

Há uma restrição prática que condiciona toda esta escolha e que convém enunciar cedo. Em domínios
financeiros vizinhos, como a deteção de fraude, os métodos \emph{não supervisionados} ocupam um lugar
central, precisamente porque anomalias rotuladas são escassas e os padrões mudam com o tempo
\autocite{ahmed2016financial}. O mesmo se aplica ao problema deste trabalho. Não existe um conjunto de dados em que alguém tenha marcado, dia a
dia, quais os movimentos que mereciam um alerta. Essa ausência não é um detalhe de implementação:
molda a forma como o Capítulo~\ref{cap:avaliacao} pode avaliar este componente, e é a razão pela
qual a avaliação se apoia num critério definido a partir dos próprios dados em vez de um julgamento
humano.

A família estatística assenta numa suposição, a de que a distribuição dos retornos é localmente
estável, e essa suposição merece ser confrontada em vez de escondida. Os modelos de volatilidade
condicional, introduzidos por \textcite{engle1982arch} sobre séries de inflação e generalizados por
\textcite{bollerslev1986garch}, formalizam o \emph{agrupamento de volatilidade}, ou seja a tendência
de os períodos calmos e os agitados chegarem em séries. Os dois artigos formalizam-na sobre séries
de inflação, e este trabalho não a toma emprestada para os retornos de ações: mede-a nos próprios
dados, na Secção~\ref{sec:av_qi1}, ao comparar a janela de pesos iguais com uma que pondera as
observações recentes com mais peso. Uma janela
deslizante é precisamente o que acompanha esse agrupamento, porque a norma contra a qual o dia de
hoje é comparado é reestimada a cada dia a partir do passado recente. A escolha não elimina o
problema, mas trata-o, e a medição referida acima diz quanto custa mantê-la.

A comparação entre a família estatística e a aprendida não fica em argumento. O
Capítulo~\ref{cap:avaliacao} corre as duas sobre os mesmos dados e reporta o resultado tal como sai.

\section{Representar texto por significado}
\label{sec:ctx_texto}

Comparar notícias por significado depende inteiramente de como se representa o texto, e essa
representação passou por três gerações, esquematizadas na Figura~\ref{fig:ctx_geracoes} e detalhadas
na Tabela~\ref{tab:ctx_texto}.

\begin{figure}[ht]
\centering
\begin{tikzpicture}[
  font=\small, >={Stealth[length=2.4mm]},
  % Altura mínima igual nas três: com alturas diferentes as caixas ficam desalinhadas e a
  % figura lê-se como se uma geração fosse maior do que as outras, que não é o que diz.
  g/.style={draw, rounded corners, align=center, text width=3.0cm, minimum height=3.3cm,
            inner sep=5pt, fill=black!3},
]
  \node[g] (a) {\textbf{Contar palavras}\\[3pt]
    \footnotesize cada palavra é\\ \footnotesize independente\\[3pt]
    \footnotesize \emph{``subiu'' e ``valorizou''}\\ \footnotesize \emph{não têm relação}};
  \node[g, right=0.85cm of a] (b) {\textbf{Palavras como vetores}\\[3pt]
    \footnotesize palavras parecidas\\ \footnotesize ficam próximas\\[3pt]
    \footnotesize \emph{mas cada palavra tem}\\ \footnotesize \emph{sempre o mesmo sentido}};
  \node[g, right=0.85cm of b] (c) {\textbf{Frases inteiras}\\[3pt]
    \footnotesize o sentido depende\\ \footnotesize do contexto\\[3pt]
    \footnotesize \emph{é o que este}\\ \footnotesize \emph{trabalho usa}};
  \draw[->, thick, black!45] (a) -- (b);
  \draw[->, thick, black!45] (b) -- (c);
\end{tikzpicture}
\caption[Três gerações de representação de texto]{Como se passou de contar palavras a representar o
sentido de uma frase inteira. Cada geração resolve uma limitação da anterior.}
\label{fig:ctx_geracoes}
\end{figure}

A primeira geração conta palavras. O modelo de espaço vetorial \autocite{salton1975vsm} representa um
documento pela frequência dos termos que contém. É transparente e é rápido, mas é cego a sinónimos:
para este tipo de representação, ``subiu'' e ``valorizou'' são símbolos sem relação nenhuma. Dentro
desta mesma geração há ainda uma segunda linha, nascida de um quadro diferente e não deste, a das
funções de ordenação probabilísticas, de que o BM25 \autocite{robertson2009bm25} é o representante
corrente; as duas continuam a ser linhas de base fortes em recuperação de informação, e partilham
aquela cegueira. Em texto financeiro há ainda uma variante
especializada, os léxicos de domínio, que corrigem o facto de dicionários de sentimento genéricos
lerem mal palavras financeiras correntes como ``passivo'' ou ``imposto''
\autocite{loughran2011liability}. São inteiramente transparentes, e estão limitados pelo vocabulário
que alguém teve o cuidado de compilar.

A segunda geração representa cada \emph{palavra} como um vetor num espaço onde a proximidade
significa uso semelhante, aprendido a partir do contexto local \autocite{mikolov2013word2vec} ou da
coocorrência global \autocite{pennington2014glove}. Isto aproxima sinónimos, que era o problema da
geração anterior, mas mantém uma limitação própria: cada palavra fica com um único vetor,
independentemente da frase em que aparece.

A terceira geração remove essa limitação. O \emph{Transformer} \autocite{vaswani2017attention} tornou
viável processar a frase inteira com dependências de longo alcance, e o BERT
\autocite{devlin2019bert} produziu representações profundamente \emph{contextuais}, em que o vetor de
uma palavra depende das que a rodeiam. Falta ainda um passo para o problema deste trabalho, porque
comparar duas frases com um modelo desses exigiria processá-las em conjunto, uma vez por par, o que
não escala para uma base com dezenas de milhares de casos. O Sentence-BERT \autocite{reimers2019sbert}
resolve isso ao produzir um vetor por frase, calculado uma única vez, de modo que a comparação passa
a ser uma operação aritmética entre vetores já existentes. É esta a representação que este trabalho
usa, e a razão da escolha é tanto de qualidade como de arquitetura.

Convém reter uma parte da contribuição do Sentence-BERT que não é a arquitetura, porque ela volta a
ser precisa no Capítulo~\ref{cap:avaliacao}. O que esse trabalho acrescenta é um \textbf{objetivo de
treino}: os vetores são aprendidos de maneira a que a distância entre dois deles signifique
semelhança de sentido. Os autores reportam que um codificador da família BERT usado diretamente para
comparar frases produz representações fracas para esse fim, ficando por vezes atrás da simples média
de vetores de palavra \autocite{reimers2019sbert}. A consequência é que um codificador afinado para
outra tarefa qualquer não herda esta propriedade só por ter sido treinado no domínio certo.

\begin{table}[ht]
\caption[Abordagens de representação de texto]{As opções consideradas para representar notícias
financeiras, e o que cada uma implica.}
\label{tab:ctx_texto}
\centering
\small
\begin{tabular}{L{0.22\textwidth} L{0.28\textwidth} L{0.36\textwidth}}
\toprule
\tabhead{Abordagem} & \tabhead{Representação} & \tabhead{Forças e limitações} \\
\midrule
Contagem de termos \newline \autocite{salton1975vsm, robertson2009bm25}
  & Frequência ponderada de palavras
  & Transparente e rápida; cega a sinónimos \\
\addlinespace
Léxicos financeiros \newline \autocite{loughran2011liability}
  & Listas de palavras com polaridade
  & Totalmente transparente; limitada ao vocabulário compilado \\
\addlinespace
Vetores de palavras \newline \autocite{mikolov2013word2vec, pennington2014glove}
  & Um vetor fixo por palavra
  & Aproxima sinónimos; um só sentido por palavra \\
\addlinespace
Frases contextuais \newline \autocite{devlin2019bert, reimers2019sbert}
  & Um vetor por frase, sensível ao contexto
  & Compara frases diretamente; modelo maior, menos legível por dentro \\
\addlinespace
Modelos de domínio \newline \autocite{liu2020finbert, huang2023finbert}
  & Como acima, ajustado a texto financeiro
  & Vocabulário do domínio; não é automaticamente melhor nesta tarefa \\
\addlinespace
Modelos grandes \newline \autocite{wu2023bloomberggpt}
  & Geração de texto sobre conhecimento amplo
  & Muito capaz; caro, opaco e fora da restrição de gratuitidade \\
\bottomrule
\end{tabular}
\end{table}

Para texto financeiro existem modelos treinados de propósito para o domínio
\autocite{liu2020finbert, huang2023finbert}. A suposição natural é que sejam melhores nesta
tarefa por isso mesmo.
Essa suposição não é automaticamente verdadeira, e o Capítulo~\ref{cap:avaliacao} mede-a em vez de a
assumir.

Este trabalho situa-se dentro da literatura mais vasta de processamento de linguagem natural
aplicado a finanças. A revisão de \textcite{kearney2014textual} cobre o estado dessa área até 2014:
os léxicos de domínio e os classificadores supervisionados sobre contagens de palavras, que eram
então os instrumentos correntes. As representações contextuais de frase, que este trabalho usa,
chegariam anos depois e não constam da revisão. Existe ainda uma linha de investigação substancial que usa texto para \emph{prever}
movimentos de preço \autocite{xing2018nlffsurvey, ding2015deep}. Este trabalho não pertence a essa
linha, e a distinção é deliberada: aqui o texto serve para encontrar casos passados comparáveis,
cujo desfecho já é conhecido e já foi medido, e não para produzir uma previsão sobre o futuro.

\section{Recuperar casos parecidos, e saber se a recuperação é boa}
\label{sec:ctx_recuperacao}

Representar cada notícia por um vetor transforma a pergunta ``que casos passados se parecem com
este?'' num problema de recuperação de informação. A medida de proximidade usada é o cosseno do
ângulo entre os vetores, cuja mecânica se detalha na Secção~\ref{sec:met_cosseno}. A operação é
exata, e portanto o resultado não depende de nenhuma aproximação: percorrem-se todos os casos da
base e ordenam-se. A escalas muito superiores, essa procura exaustiva pode ser trocada por um índice
aproximado \autocite{johnson2021faiss} sem alterar a representação, o que fica registado como
caminho de crescimento e não como necessidade atual.

Como a recuperação devolve uma lista ordenada, avalia-se com medidas sensíveis à ordem
\autocite{manning2008ir}. A mais direta, e a que este trabalho usa, é a \emph{precisão@$k$}, ou seja
a fração dos $k$ primeiros resultados que são relevantes. A definição de relevância, que é onde uma
avaliação destas se ganha ou se perde, é estabelecida no Capítulo~\ref{cap:metodologia} juntamente com as
linhas de base contra as quais o resultado é comparado.

\section{Medir o efeito de um acontecimento}
\label{sec:ctx_evento}

O instrumento é o \emph{estudo de evento}, introduzido por \textcite{fama1969adjustment}.
\textcite{brown1985daily} estudam em detalhe o seu comportamento com retornos \emph{diários},
comparando várias formas de estimar o retorno esperado, e concluem que os procedimentos correntes,
assentes no modelo de mercado estimado por mínimos quadrados, se mantêm bem especificados em janelas
curtas em torno do acontecimento. Os próprios autores delimitam o alcance dessa conclusão: a variante
mais simples de todas, que desconta apenas a média histórica da própria ação, perde potência e
chega a ficar mal especificada quando as datas dos acontecimentos se concentram no mesmo dia.

É esse resultado que dá conforto às janelas de mais um, mais três e mais cinco dias usadas neste
trabalho, com uma diferença que a Secção~\ref{sec:met_recuperacao} declara: o que aqui se mostra ao
utilizador é o retorno bruto e não o retorno anormal que aqueles autores estudam.

Que o texto publicado tem relação estatística com os movimentos do mercado é observação antiga
\autocite{tetlock2007media}. Convém enunciá-la com o verbo certo: trata-se de evidência de uma
relação estatística, e não de uma prova de causalidade.

Há um ponto metodológico que interessa antecipar aqui, porque afeta a segunda pergunta do trabalho,
a de separar a empresa do mercado. Fazer essa separação exige estimar a sensibilidade de cada ação
ao mercado, e essa estimativa é ruidosa, sobretudo em empresas voláteis ou com histórico curto. A
literatura de finanças documentou o fenómeno e propôs a correção. \textcite{blume1971risk} estabelece
que os betas estimados tendem a regredir para a média entre períodos sucessivos, ou seja que uma
estimativa extrema hoje tende a sê-lo menos amanhã; daí decorreu a prática, corrente mas posterior ao
próprio artigo, de encolher a estimativa com um peso fixo. Essa prática encolhe com a mesma força uma
estimativa boa e uma má. A alternativa que este trabalho adota pondera o encolhimento pela
\emph{imprecisão} da própria estimativa \autocite{vasicek1973beta}, de modo que uma estimativa fiável
é praticamente preservada e uma frágil é puxada para o valor de referência. O detalhe está na
Secção~\ref{sec:met_decomposicao}.

Por fim, o enquadramento teórico que justifica a recusa de prever. A hipótese dos mercados
eficientes, na sua forma semiforte, sustenta que os preços já refletem a informação publicamente
disponível \autocite{fama1970efficient}; daí decorre que movimentos futuros não deveriam ser
sistematicamente previsíveis a partir dessa informação. Esta é uma das razões pelas quais o sistema mede o que já
aconteceu em vez de projetar o que virá, e a outra, mais prática, é que apenas o passado é
verificável por quem lê.

\section{Raciocinar por casos passados}
\label{sec:ctx_casos}

Tem, e é uma linha estabelecida em inteligência artificial. O \emph{raciocínio baseado em casos}
\autocite{aamodt1994cbr} resolve um problema novo recuperando problemas passados semelhantes e
reutilizando o que se sabe sobre eles, em vez de se apoiar \emph{apenas} em conhecimento geral do
domínio. A fonte descreve um ciclo de quatro processos, que são recuperar, reutilizar, rever e reter,
e afirma que os quatro são necessários para resolver um problema por completo.

A terceira técnica deste trabalho implementa os dois primeiros e \textbf{para deliberadamente} antes
dos outros dois: recupera casos semelhantes e mostra o que neles foi medido, mas não adapta esse
desfecho ao caso presente. Adaptar seria produzir uma expectativa sobre o que vai acontecer, ou seja
exatamente a previsão que este trabalho recusa.

A escolha tem uma justificação adicional, que vem do lado de quem lê. A investigação sobre explicação
em ciências sociais mostra que as pessoas procuram razões \emph{contrastivas} e seletivas, e não uma
enumeração exaustiva de tudo o que contribuiu \autocite{miller2019explanation}. Um punhado de casos
concretos, com datas e desfechos, corresponde a essa expectativa muito melhor do que um coeficiente.
Daí que cada alerta mostre alguns precedentes relevantes e as quantidades exatas por detrás da
decisão, em vez de despejar tudo o que o sistema calculou.

\section{Explicar decisões automáticas}
\label{sec:ctx_xai}

A exigência de que um sistema mostre \emph{como} chegou a uma conclusão constitui uma área própria,
com revisões consolidadas \autocite{arrieta2020xai, adadi2018peeking}. Essa área divide-se em duas
vias, ilustradas na Figura~\ref{fig:ctx_xai}: modelos
\emph{transparentes}, interpretáveis em si mesmos, e métodos \emph{a posteriori}, que explicam um
modelo opaco já treinado e que podem ser organizados pelo que explicam
\autocite{guidotti2018survey}.

\begin{figure}[ht]
\centering
\begin{tikzpicture}[
  font=\small, >={Stealth[length=2.4mm]},
  r/.style={draw, rounded corners, align=center, text width=3.5cm, minimum height=1.1cm,
            inner sep=4pt, fill=black!3},
  t/.style={draw, rounded corners, align=center, text width=3.5cm, minimum height=1.4cm,
            inner sep=4pt, fill=black!8, very thick},
  o/.style={draw, rounded corners, align=center, text width=3.5cm, minimum height=1.4cm,
            inner sep=4pt, fill=black!3},
]
  \node[r] (raiz) {\textbf{Explicar uma decisão}};
  \node[t, below left=1.1cm and 0.35cm of raiz] (tr) {\textbf{Modelo transparente}\\
    \footnotesize a explicação é o cálculo};
  \node[o, below right=1.1cm and 0.35cm of raiz] (ph) {\textbf{Método a posteriori}\\
    \footnotesize explica um modelo opaco};

  \node[below=0.4cm of tr, align=center, font=\footnotesize, text width=3.6cm]
    (trx) {regra, regressão linear,\\ contagem sobre a série};
  \node[below=0.4cm of ph, align=center, font=\footnotesize, text width=3.6cm]
    (phx) {LIME, SHAP};

  \draw[->, thick, black!45] (raiz) -- (tr);
  \draw[->, thick, black!45] (raiz) -- (ph);
  \draw[-, black!35] (tr) -- (trx);
  \draw[-, black!35] (ph) -- (phx);
\end{tikzpicture}
\caption[As duas vias para a explicabilidade]{As duas vias que a literatura distingue
\autocite{arrieta2020xai, guidotti2018survey}: construir modelos transparentes por desenho, ou
explicar um modelo opaco depois de treinado. Este trabalho segue a via da esquerda.}
\label{fig:ctx_xai}
\end{figure}

Os dois métodos \emph{a posteriori} mais usados operam de formas distintas. O \gls{LIME} ajusta um modelo
simples em torno de uma previsão individual, aproximando localmente o comportamento do modelo opaco
\autocite{ribeiro2016lime}. O \gls{SHAP} atribui uma previsão às suas variáveis usando valores de Shapley
\autocite{lundberg2017shap}. Ambos são genéricos e ambos são úteis.

Ambos partilham, contudo, uma propriedade que pesa muito neste contexto. Para decisões de
consequência séria, \textcite{rudin2019stop} argumenta que se deve abandonar o compromisso com
modelos opacos e usar antes modelos interpretáveis à partida, precisamente porque uma explicação
\emph{a posteriori} é ela própria uma aproximação, e uma aproximação pode induzir em erro. O
utilizador recebe então duas coisas cuja relação não consegue avaliar: uma decisão e uma explicação
que talvez não corresponda a ela.


Este trabalho segue a via conservadora: em vez de explicar \emph{a posteriori} um modelo opaco, usa
componentes cuja explicação é o próprio cálculo. Um desvio relativo a uma norma recente, uma
decomposição aditiva que soma exatamente ao movimento observado e uma lista de casos passados com os
respetivos desfechos são todos objetos que se podem mostrar por inteiro. A explicação não é uma
segunda coisa produzida sobre a decisão: é a decisão, escrita por extenso.

Duas cautelas da literatura fecham esta secção, e as duas condicionam o modo como o
Capítulo~\ref{cap:avaliacao} está organizado.

A primeira é que ``interpretabilidade'' não designa uma coisa só \autocite{lipton2018mythos}, pelo
que um sistema deve declarar que tipo de interpretabilidade oferece em vez de reclamar a palavra sem
qualificação. Este trabalho oferece transparência de mecanismo, ou seja o utilizador pode seguir o
cálculo, e não oferece uma explicação causal do acontecimento no mundo.

A segunda é mais incómoda e é frequentemente ignorada. Explicações podem induzir
\emph{confiança excessiva}: \textcite{bansal2021whole} mostram que acrescentar uma explicação
aumenta a probabilidade de a pessoa aceitar a resposta do sistema \textbf{quer ela esteja certa quer
esteja errada}, ou seja que o ganho não vem de a pessoa passar a distinguir melhor os dois casos. É
a presença da explicação que produz o efeito, e não o tom com que é apresentada. Daqui decorre uma
consequência de desenho concreta neste sistema. A média dos desfechos dos casos recuperados é
mostrada como aquilo que é, uma medição sobre o passado, e nunca como uma expectativa sobre o futuro,
porque tratá-la como previsão seria exatamente a confiança excessiva contra a qual a literatura de
confiança avisa \autocite{lee2004trust, bansal2021whole}.

Acresce que afirmações de interpretabilidade exigem avaliação rigorosa e não basta declará-las
\autocite{doshivelez2017rigorous}. É por isso que o Capítulo~\ref{cap:avaliacao} mede o que pode ser
medido por máquina e diz, de forma explícita, qual a parte que continua por avaliar por faltar um
estudo com pessoas.

\section{Um modelo que continua a funcionar depois de implantado}
\label{sec:ctx_producao}

O componente aprendido deste trabalho levanta três questões que a literatura trata separadamente.

A primeira é contra que teto comparar um modelo simples. Conjuntos de árvores treinados por reforço
do gradiente \autocite{friedman2001gbm} captam interações e efeitos não lineares que uma regressão
logística não capta, e servem por isso de referência superior no Capítulo~\ref{cap:avaliacao}. Não
se afirma que sejam o melhor método existente, porque essa seria uma afirmação sobre um campo
inteiro.

A segunda é a honestidade do número produzido. Uma pontuação só é honesta para quem a lê se for uma
probabilidade, e as pontuações em bruto de muitos classificadores estão mal calibradas; a
calibração \emph{a posteriori} sobre um bloco de validação reservado corrige essa distorção
\autocite{niculescu2005calibration}. A mesma fonte, porém, não favorece a escolha feita aqui:
aponta a regressão logística como já bem calibrada à partida, e conclui que acima de cerca de mil
casos de calibração a regressão isotónica iguala ou supera a sigmoide. O modelo implantado é uma
regressão logística com um bloco de validação muito acima desse limiar, pelo que ambas as
afirmações se aplicam. A calibração justifica-se aqui por medição e não por autoridade: a
Secção~\ref{sec:av_alternativas} compara as duas variantes na mesma validação e a sigmoide iguala
ou ganha no Brier em todas as famílias testadas.

Uma via alternativa foi considerada e não adotada. A predição conformal devolve um \emph{conjunto}
de rótulos com garantia livre de distribuição válida em amostras finitas
\autocite{vovk2005algorithmic, angelopoulos2023conformal}. Duas precisões decidem que não serve
aqui: a cobertura da variante corrente é \emph{marginal}, vale em média sobre os casos e não caso a
caso; e a garantia exige \textbf{permutabilidade} entre calibração e teste, que é exatamente a
hipótese que a Secção~\ref{sec:met_triagem} recusa ao dividir o conjunto por data e com embargo. O
Capítulo~\ref{cap:conclusoes} retoma o método como caminho a explorar.

A terceira questão é continuar a funcionar com o tempo. Um modelo implantado enfrenta alteração da
distribuição dos dados que vê, e a literatura de \emph{deriva de conceito} estabelece a taxonomia
da mudança e as estratégias de deteção e adaptação \autocite{gama2014survey}. Este trabalho mede a
deriva sobre os seus próprios dados em vez de a assumir ausente.

Há por fim um custo que não aparece em nenhuma métrica de qualidade. \textcite{sculley2015debt}
argumentam que os componentes de aprendizagem automática acumulam manutenção oculta, em particular
o emaranhamento entre peças e as dependências de dados instáveis e não declaradas. O
Capítulo~\ref{cap:conclusoes} regressa a este ponto com um exemplo medido deste sistema, em que
parte do ciclo de aprendizagem parou sem que nenhum erro fosse registado.

\section{Síntese: a lacuna que este trabalho ocupa}
\label{sec:ctx_lacuna}

A revisão feita neste capítulo permite localizar com precisão o espaço em que este trabalho se
inscreve, e esse espaço não é um vazio técnico. Cada peça individual existe, está publicada e está
avaliada. A Tabela~\ref{tab:ctx_lacuna} torna isso explícito.

\begin{table}[ht]
\caption[Onde fica este trabalho]{Cada componente existe na literatura. O que não existe é a
combinação, sob a restrição de custo nulo, com a evidência entregue ao utilizador.}
\label{tab:ctx_lacuna}
\centering
\small
\begin{tabular}{L{0.24\textwidth} L{0.30\textwidth} L{0.34\textwidth}}
\toprule
\tabhead{Pergunta} & \tabhead{O que a literatura oferece} & \tabhead{O que falta para esta pessoa} \\
\midrule
É invulgar?
  & Deteção de anomalias, estatística e aprendida \autocite{chandola2009anomaly, liu2008isolation}
  & Chegar-lhe em linguagem que ela leia, relativa à norma daquela empresa \\
\addlinespace
Empresa ou mercado?
  & Estudo de evento e estimação de sensibilidade \autocite{brown1985daily, vasicek1973beta}
  & Estar disponível fora de um terminal profissional \\
\addlinespace
Já aconteceu antes?
  & Representação de frases e recuperação \autocite{reimers2019sbert, manning2008ir}
  & Uma base de casos com desfechos já medidos, e não recordados \\
\addlinespace
Porque é que acredito?
  & Explicabilidade e literatura de confiança \autocite{rudin2019stop, lee2004trust}
  & A evidência entregue junto da afirmação, verificável no momento da leitura \\
\bottomrule
\end{tabular}
\end{table}

A lacuna, portanto, não é de método: é de \emph{integração sob restrição}. Os produtos que respondem
às três perguntas cobram milhares por ano. Os produtos gratuitos que prometem responder à terceira
entregam uma afirmação sem a evidência. E os componentes académicos que resolveriam cada pedaço
foram avaliados isoladamente, sobre corpora de referência, e não como um sistema que funciona todos
os dias contra fontes reais e que reporta o que falha.

É esse o espaço que este trabalho ocupa, e é também o que determina a natureza da sua contribuição:
não está em inventar um método novo, mas em escolher, integrar e avaliar criticamente métodos
existentes num sistema que funciona, sob restrições declaradas, e cuja avaliação inclui os
resultados que não favorecem o próprio sistema.

\section{As escolhas que este trabalho faz}
\label{sec:ctx_escolhas}

Quatro decisões, e cada uma fecha uma porta de propósito.

\begin{description}
  \item[Responder às três perguntas, e só a essas.] Não se acrescentam funcionalidades que não
        sirvam uma delas. É o que mantém o sistema pequeno o suficiente para ser explicado por
        inteiro.
  \item[Não prever preços.] Nem direção, nem alvo, nem recomendação. Além da razão já dada, que é o
        que torna tudo verificável, há a razão teórica da Secção~\ref{sec:ctx_evento}: os preços já
        refletem a informação publicamente disponível \autocite{fama1970efficient}, e se assim for,
        movimentos futuros não deveriam ser sistematicamente previsíveis a partir dela.
  \item[Não gerir carteiras nem aconselhar.] O sistema não sabe o que o utilizador tem, e não lhe diz
        o que fazer. Isto evita recolher dados pessoais e mantém o trabalho fora da fronteira do
        aconselhamento financeiro regulado.
  \item[Só fontes gratuitas.] Torna o sistema utilizável por quem ele serve, e transforma as
        limitações resultantes em algo que se \emph{mede} e reporta, em vez de algo que se compra.
\end{description}

A última merece uma nota honesta. A restrição limita o sistema de forma real: as notícias chegam
com atraso, e nem tudo o que acontece é coberto. O Capítulo~\ref{cap:sistema} quantifica esse
atraso e o Capítulo~\ref{cap:conclusoes} volta a ele. É uma limitação, e não uma virtude.
